\documentclass[10pt,a4paper]{amsart}
\usepackage[margin=19mm]{geometry}
\usepackage{amsmath,amssymb,mathtools}
\usepackage[scr=rsfs,cal=euler]{mathalfa}
\usepackage{xfrac}
\usepackage[unicode,hidelinks,bookmarks=false]{hyperref}
\newcommand{\ie}{i.e.\ }
\newcommand{\cf}{cf.\ }
\newcommand{\resp}{respectively\ }
\newcommand{\Subsection}{\subsection}
\providecommand{\nobreakdash}{\nobreak\hbox{-}}
\setlength{\emergencystretch}{2em}
\multlinegap0pt
\usepackage{bm}
\usepackage{bbm}
\usepackage{dsfont}
\usepackage{xspace}
\usepackage{cancel}
\usepackage{mathtools}
\usepackage{amsthm}
\makeatletter
\@ifundefined{theorem}{\newtheorem{theorem}{Theorem}}{}
\@ifundefined{definition}{\newtheorem{definition}[theorem]{Definition}}{}
\@ifundefined{lemma}{\newtheorem{lemma}[theorem]{Lemma}}{}
\@ifundefined{proposition}{\newtheorem{proposition}[theorem]{Proposition}}{}
\makeatother
\newcommand{\A}{\mathcal{A}}
\newcommand{\PP}{\mathbb{P}}
\newcommand{\cube}{\odot^3}
\title[Abelian surfaces in Hesse form and explicit isogeny formulas]{Abelian surfaces in Hesse form and explicit isogeny formulas}
\author{Thomas Decru, Sabrina Kunzweiler}
\date{Source: arXiv:2601.05922v1 (CC BY 4.0)}
\begin{document}
\maketitle

\noindent\emph{Source and licence.} Abelian surfaces in Hesse form and explicit isogeny formulas. Version arXiv:2601.05922v1 (CC BY 4.0), distributed under the Creative Commons Attribution 4.0 International licence, as recorded in the arXiv licence metadata for this exact version and shown on the abstract page https://arxiv.org/abs/2601.05922v1. Full rights notice: licenses/arxiv-2601.05922-CC-BY-4.0.txt.\par\smallskip

\noindent\textit{Editorial scope.} Counted unit: the Proposition whose source label is \texttt{prop:isogeny-formula-exceptional} in the article (source0.tex lines 1242--1258, source label \texttt{prop:isogeny-formula-exceptional}), delivered together with the article's own complete original proof of it (source0.tex lines 1260--1312, one \texttt{proof} environment of 53 lines). No mathematics is written, repaired or re-derived: statement and proof are the article's own text line for line. Required context retained uncounted: def:exceptional (lines 1080--1090); thm:3-isogeny-dim2 (lines 1093--1117); lem:exceptional-points (lines 1233--1235). Every definition, display and label of those lines is retained unchanged. Omitted from this package and not redistributed: 1--1079, 1091--1092, 1118--1232, 1236--1241, 1259--1259, 1313--1529. Nothing inside a retained excerpt is omitted or altered: every retained body is delivered exactly as the article's lines are written, so no omission receipt of any kind accompanies this package. Citations: the retained excerpts carry no citation command at all, so no bibliography entry of the article is transcribed and no reference is redistributed. Reference adaptation: none was needed and none was made; every reference inside the delivered unit resolves inside it, so no unresolved reference or citation is produced. Printed slips kept literal: no printed word, symbol, hypothesis or number of the counted statement or of its proof was altered. Layout and class: the article's own class is \texttt{amsart} with packages including inputenc, xcolor, amsmath, amssymb, amsfonts, hyperref, mathtools, tikz-cd, array, orcidlink, amsthm, placeins, fullpage, multicol; the wrapper is the lane's pinned \texttt{amsart} editorial wrapper at 10pt on A4, which restates the theorem-like environments the retained text needs and adds the article's own macro definitions that the retained text needs (\texttt{\textbackslash A}, \texttt{\textbackslash PP}, \texttt{\textbackslash cube}), with extra wrapper packages (bm, bbm, dsfont, xspace, cancel, mathtools). The delivered numbering is the unit's own. Assets: no retained span contains a figure environment, a graphics-inclusion command, a PDF-page inclusion, a table or a TikZ picture; no image file is copied into this package, the package declares no asset dependency and no image file is redistributed with it. Licence: version arXiv:2601.05922v1 of this article is distributed under the Creative Commons Attribution 4.0 International licence, as recorded in the arXiv licence metadata for this exact version and shown on the abstract page https://arxiv.org/abs/2601.05922v1. A full rights notice is delivered at licenses/arxiv-2601.05922-CC-BY-4.0.txt. Characters: every retained excerpt is pure ASCII, so the delivered engine, XeLaTeX with its default Latin Modern text font, reports no missing character.
\begin{definition}\label{def:exceptional}
    Let $\A_{d,h}$ be a p.p. abelian surface in Hesse form, and $0_{\A_{d,h}} = (t_0: \dots : t_3) \in \PP^8$ the neutral element, and define the following values 
    \begin{align*}
    \tilde{t}_0 =&~ t_0^3 + 2t_1^3 + 2t_2^3 + 2t_3^3 + 2t_4^3,\\
    \tilde{t}_1 =&~  t_0^3 - t_1^3 + 2t_2^3 - t_3^3 -t_4^3,\\
    \tilde{t}_2 = &~ t_0^3 + 2  t_1^3 - t_2^3 - t_3^3 -t_4^3,\\
    \tilde{t}_3 = &~ t_0^3 - t_1^3 - t_2^3 - t_3^3  +2 t_4^3,\\
    \tilde{t}_4 = &~ t_0^3 - t_1^3 - t_2^3 +2 t_3^3 - t_4^3.
    \end{align*}
    We say that $\A_{d,h}$ is {\em exceptional} if there is an index $i>0$ with $\tilde{t}_i = 0$.
\end{definition}

\begin{theorem} \label{thm:3-isogeny-dim2}
    Let $\A_{d,h}$ be a p.p. abelian surface in Hesse form, and assume that $\A_{d,h}$ is not exceptional. Denote $(P_1,P_2,Q_1,Q_2)$ for the canonical $3$-torsion basis, and let 
    \(
    \langle R_1, \; R_2\rangle \subset \A_{d,h}[9]
    \) be a maximal isotropic subgroup with $3 \cdot R_1  =  Q_1$ and $3 \cdot R_2  =  Q_2$. We denote
    \[
    (x_{i,0}:\dots: x_{i,8}) = M \circ \cube(R_i) \quad \text{for } i =1,2.
    \]
    and set 
    \begin{align*}
    (\lambda_0 : \lambda_1) =&~ (x_{1,1} :  x_{1,0}),
    \\
     (\lambda_0 : \lambda_2) =&~ (x_{2,3} :  x_{2,0}),
     \\
    (\lambda_2 : \lambda_3) =&~ (x_{1,4}: x_{1,6}),
    \\ 
    (\lambda_3 : \lambda_4) =&~ (x_{1,5} : x_{1,8}).  
    \end{align*}
    
    Then 
    \[
    \phi: \A_{d,h} \to \A_{d',h'}, \quad P \mapsto C_\lambda \circ M \circ \cube(P)
    \]
    defines a $(3,3)$-isogeny with kernel $K_2 = \langle Q_1, Q_2\rangle$ such that the codomain is again in Hesse form. Moreover, the canonical symplectic basis for $\A_{d',h'}[3]$ is given by  $(P_1',P_2',Q_1',Q_2')$, where $P_1' = \phi(R_1), \, P_2' = \phi(R_2)$ and $Q_1' = [-1]\phi(P_1), \, Q_2' = [-1]\phi(P_2)$. 
\end{theorem}

\begin{lemma} \label{lem:exceptional-points}
    Let $\A_{d,h}$ be an exceptional p.p. abelian surface. Let $\tilde{t}=(\tilde t_0: \dots: \tilde t_4) $ as in  Definition \ref{def:exceptional}. Then there exist at least two indices $i \in \{0, \dots, 4\}$  with $\tilde t_i \neq 0$.
\end{lemma}

\begin{proposition} \label{prop:isogeny-formula-exceptional}
    Let $\A_{d,h}$ be a p.p. abelian surface in Hesse form. Denote $(P_1,P_2,Q_1,Q_2)$ for the canonical $3$-torsion basis, and let 
    \(
    \langle R_1,R_2\rangle \subset \A_{d,h}[9]
    \) be a maximal isotropic subgroup with $3 \cdot R_1  =  Q_1$, $3 \cdot R_2  =  Q_2$  and $R_3 = R_1 + R_2$, $R_4 = R_1 - R_2$. %We denote
    %\[
    %(x_{i,0}:\dots: x_{i,8}) = M \circ \cube(R_i) \quad \text{for } i =1,2,3,4.
    %\]
    
    Further, let $\tilde{t} = (\tilde{t}_0: \dots: \tilde{t}_4)$ as in Definition \ref{def:exceptional}. If there are two indices $i \neq j >0$ with $\tilde{t}_i, \tilde{t}_j \neq 0$, then there exists a unique $\lambda \in \PP^4$ so that 
    \[
    \phi: \A_{d,h} \to \A_{d',h'}, \quad P \mapsto C_\lambda \circ M \circ \cube(P)
    \]
    defines a $(3,3)$-isogeny with kernel $K_2 = \langle Q_1, Q_2\rangle$, and the canonical symplectic basis for $\A_{d',h'}[3]$ is given by  $(P_1',P_2',Q_1',Q_2')$, where $P_1' = \phi(R_1), \, P_2' = \phi(R_2)$ and $Q_1' = [-1]\phi(P_1), \, Q_2' = [-1]\phi(P_2)$. 
    
    Moreover, the coordinates of $\lambda$ can be read off from the $9$-torsion points $R_1,R_2,R_3,R_4$.
\end{proposition}

\begin{proof}
	We use the same notation as in the proof of Theorem \ref{thm:3-isogeny-dim2}. That is we decompose the isogeny $A \xrightarrow{\phi_2} A_2 \xrightarrow{C_\lambda} A_{d',h'}$ for some $\lambda \in \PP^4$. Let $(P_1',P_2',Q_1',Q_2')$ be the canonical symplectic basis of $\A_{d',h'}$. With the same arguments as in the proof of the theorem, we may assume that $\phi(R_1) = P_1'$ and $\phi(R_2)=P_2'$, hence $\phi(R_3) = P_1'+P_2' =: P_3'$ and $\phi(R_4) = P_1'-P_2' =:P_4'$.
	
	
	We denote
	\[
	\phi_2(R_i) = (x_{i,0}:\dots: x_{i,8}) \quad \text{for } i =1,2,3,4.
	\]
	The condition $C_\lambda(\phi_2(R_i)) = P_i'$ for $i = 1,2,3,4$ provides us with four equalities. 
	
	\begin{align*}
		\left(\lambda_0 x_{1,0}: \lambda_1 x_{1,1} : \dots: \lambda_3 x_{1,8}\right) = &~
		(t_1' : t_1' : t_0':  t_3': t_4' : t_2' :  t_4' : t_3' : t_2'),\\
		\left(\lambda_0 x_{2,0}: \lambda_1 x_{2,1} : \dots: \lambda_3 x_{2,8}\right)  = &~
		(t_2': t_3' : t_4' : t_2' : t_4' : t_3 ': t_0' : t_1' : t_1' ),\\
		\left(\lambda_0 x_{3,0}: \lambda_1 x_{3,1} : \dots: \lambda_3 x_{3,8}\right)  = &~
		(t_3': t_4' : t_2' : t_4' : t_3' : t_2 ': t_1' : t_1' : t_0' ),\\
		\left(\lambda_0 x_{4,0}: \lambda_1 x_{4,1} : \dots: \lambda_3 x_{4,8}\right)  = &~
		(t_4': t_3' : t_2' : t_1' : t_1' : t_0 ': t_3' : t_4' : t_2' ).
	\end{align*}
	
	For each $i$, the condition $\tilde{t_i} \neq0$ is equivalent to $t_i' \neq 0$ (cf. the proof of Lemma \ref{lem:exceptional-points}). Now if $t_i'\neq 0$ for some $i>0$, then we can deduce one relation on the coordinates of $\lambda$ from each of the above equalities.\\
	Explicitly, if $t_1' \neq 0$, then
	\[
	(\lambda_0:\lambda_1) = (x_{1,1}:x_{1,0}), ~~
	(\lambda_4:\lambda_3) = (x_{2,8}:x_{2,7}), ~~
	(\lambda_2:\lambda_4) = (x_{3,7}:x_{3,6}), ~~
	(\lambda_2:\lambda_3) = (x_{4,4} : x_{4,3}).
	\]
	If $t_2' \neq 0$, then
	\[
	(\lambda_4:\lambda_3) = (x_{1,8}:x_{1,5}), ~~
	(\lambda_0:\lambda_2) = (x_{2,3}:x_{2,0}), ~~
	(\lambda_1:\lambda_4) = (x_{3,5}:x_{3,2}), ~~
	(\lambda_1:\lambda_3) = (x_{4,8} : x_{4,2}).
	\]
	If $t_3' \neq 0$, then
	\[
	(\lambda_2:\lambda_4) = (x_{1,7}:x_{1,3}), ~~
	(\lambda_1:\lambda_4) = (x_{2,5}:x_{2,1}), ~~
	(\lambda_0:\lambda_3) = (x_{3,4}:x_{3,0}), ~~
	(\lambda_1:\lambda_2) = (x_{4,6} : x_{4,1}).
	\]
	If $t_4' \neq 0$, then
	\[
	(\lambda_3:\lambda_2) = (x_{1,6}:x_{1,4}), ~~
	(\lambda_1:\lambda_3) = (x_{2,4}:x_{2,2}), ~~
	(\lambda_1:\lambda_2) = (x_{3,3}:x_{3,1}), ~~
	(\lambda_0:\lambda_4) = (x_{4,7} : x_{4,0}).
	\]
	
	Note that in all of the above cases, there is redundancy in the derived relations. In particular, using only one set of relations obtained for some $t_i' \neq 0$ is not enough to recover the coordinate vector $\lambda \in \PP^4$. However, two sets of relations combined for some $t_i',t_j' \neq 0$ with $i\neq j > 0$ is enough to recover $\lambda$. 
\end{proof}

\end{document}
